% This file is part of the modular source for Stratified Relational Workspace Theory.
% !TeX root = ../Stratified_Relational_Workspace_Theory.tex
\theorypart{XIX}{Summary of Main Results}

\topic{68}{Structural and Foundational Results}

\subtopic{\(T_1\): Limit--holonomy theorem}
\[
\Gamma(H)\cong C_o^{\operatorname{Hol}_o}.
\]

\subtopic{\(T_2\): Operator closure theorem}
\[
H_e\rho_i(a)=\rho_j(a)H_e
\Longrightarrow
\rho_o(a)C_*\subseteq C_*.
\]

\subtopic{\(T_{2a}\): Workspace commutant theorem }
\[
\mathcal O_{\Workspace}
\cong
\operatorname{Comm}(\Hol_o).
\]

\subtopic{\(T_3\): Transfer-Laplacian kernel theorem}
\[
\ker L_H=\Gamma(H).
\]

\subtopic{\(T_4\): Spectral stability theorem}
If \(\eta=\|\widetilde L_H-L_H\|<g_d/2\), then
\[
\|\widetilde P_d-P_d\|
\leq
\frac{\eta}{g_d-\eta}
\leq
\frac{2\eta}{g_d}.
\]

\subtopic{\(T_5\): Gauge invariance theorem}
\[
\Gamma(H^g)=G_V\Gamma(H),
\]
and under unitary gauge,
\[
L_{H^g}=G_VL_HG_V^{-1}.
\]

\subtopic{\(T_6\): Persistence theorem}

For the nonautonomous recurrence
\[
e_{t+1}\leq\alpha_te_t+\sigma_t,
\]
\[
\boxed{
\begin{aligned}
e_t
\leq
&\left(\prod_{r=0}^{t-1}\alpha_r\right)e_0\\
&+\sum_{j=0}^{t-1}
\left(\prod_{r=j+1}^{t-1}\alpha_r\right)\sigma_j.
\end{aligned}
}
\]
If \(\alpha_t\leq\bar\alpha<1\) and \(\sigma_t\leq\bar\sigma\), then
\[
\limsup_{t\to\infty}e_t
\leq
\frac{\bar\sigma}{1-\bar\alpha}.
\]
The derivation begins with a moving workspace projector and
an explicit system update, and adds a mean-square analogue and a non-collapse
gate.

\subtopic{\(T_7\): Stable reconstruction theorem}

For coalition analysis operator \(A_S\) with lower frame bound
\(\alpha_S>0\), observations
\[
y=A_Sc+\eta
\]
and canonical decoder \(A_S^\dagger\) satisfy
\[
\boxed{
\|A_S^\dagger y-c\|
\leq
\frac{\|\eta\|}{\sqrt{\alpha_S}}.
}
\]
The theorem is extended to encoding mismatch, exact support-based erasure
distance, uniform node erasures, reconstructive handoff, exact coalition
sufficiency, and operational risk under a Lipschitz predictive kernel.

\subtopic{\(T_8\): Exact gluing theorem}

For every finite edge-complete workspace cover,
\[
\boxed{
0
\longrightarrow
\mathscr W(G)
\xrightarrow{\ R_{\mathcal U}\ }
\check C^0(\mathcal U;\mathscr W)
\xrightarrow{\ \delta^0\ }
\check C^1(\mathcal U;\mathscr W)
}
\]
is exact. Hence
\[
\boxed{
\mathscr W(G)
\cong
\ker\delta^0
\cong
\operatorname{Eq}
\left(
\prod_\alpha\mathscr W_\alpha
\rightrightarrows
\prod_{\alpha<\beta}\mathscr W_{\alpha\beta}
\right).
}
\]
The same analysis also yields a cover-Laplacian distance bound, a total descent kernel
theorem, and a harmonic \(\check H^1\) obstruction for relative-only data.

\subtopic{\(T_9\): Hierarchical collapse theorem}

For every finite rooted hierarchy whose children edge-cover their parent,
\[
\boxed{
W_a\cong\mathscr W(G_a)
\quad\text{for every hierarchy node }a,
}
\]
and in particular
\[
\boxed{W_{\mathrm{root}}\cong\mathscr W(G).}
\]
If level \(r\) contributes defect \(\delta_r\) and downstream collapse maps
have Lipschitz constants \(L_s\), then
\[
\boxed{
\delta_{\mathrm{flat}}^{(R)}
\leq
\sum_{r=0}^{R-1}
\left(
\prod_{s=r+1}^{R-1}L_s
\right)\delta_r.
}
\]


\subtopic{\(T_{25}\): Faithful coarse-graining theorem}

For a partition of the vertex set, let \(W_\Pi\) be the direct sum of the
internally coherent section spaces and let \(D_\Pi\) retain all boundary
compatibility equations. Then
\[
\boxed{
\Gamma(H)
=
J_\Pi\ker D_\Pi.
}
\]
Exact structural workspace content is therefore invariant under faithful
aggregation. Approximate error separates into an internal-coherence term and
a boundary-coherence term controlled by their respective spectral gaps.

\subtopic{\(T_{26}\): Bayes-risk transport theorem}

For an information-faithful coarsening \(q:V_f\to V_c\),
\[
\boxed{
\Risk_{\mathcal D_c}(S)
=
\Risk_{\mathcal D_f}(q^{-1}(S))
\qquad(S\subseteq V_c).
}
\]
Hence displayed coalition size can decrease under coarsening. Weighted
primitive support cost is a more stable localization measure and cannot
improve merely because several carriers were relabeled as one macro-node.

\subtopic{\(T_{27}\): Decomposition-transport robustness theorem}

If every gate margin at a reference decomposition exceeds its worst-case
transport and model-mismatch defect across \(\mathfrak D\), then the complete
certificate passes throughout the admissible family:
\[
\boxed{
 m_q(W;\mathcal D_0)
>
\sup_{\mathcal D\in\mathfrak D}
\bigl(L_q\delta_{\mathcal D}+b_q(\mathcal D)\bigr)
\quad\forall q
\Longrightarrow
\operatorname{Certify}_{L,\mathfrak D}^{\mathrm{rob}}(W)=1.
}
\]

\topic{69}{Dynamic, Causal, and Agency Results}

\subtopic{\(T_{10}\): Approximate universality theorem }

For the first \(d\) transfer-Laplacian modes, with
\(\lambda_d\leq\varepsilon<\lambda_{d+1}=\mu\),
\[
\boxed{
\|(I-P_d)\rho_V(a)P_d\|
\leq
\frac{\|\rho_E(a)\|\sqrt{\varepsilon}+\beta(a)}{\sqrt{\mu}}.
}
\]
Thus low transfer energy, a separated complement, and small intertwining defect
imply approximate closure under the common operation algebra.

\subtopic{\(T_{11}\): Core--residual accumulation theorem }

For a path \(\gamma=e_m\cdots e_1\),
\[
\boxed{
\begin{aligned}
\|Q_{t(\gamma)}T_\gamma P_{s(\gamma)}\|
\leq
\sum_{r=1}^{m}
&\left(
\prod_{\ell=r+1}^{m}\|T_{e_\ell}\|
\right)
\|V_{e_r}\|\\
&\times
\left(
\prod_{\ell=1}^{r-1}\|H_{e_\ell}\|
\right).
\end{aligned}
}
\]
The same bound controls the defect between full core-to-core path transport and
the pure core product. An exact Duhamel
identity yields both estimates.

\subtopic{\(T_{12}\): Predictive-state minimality and alignment theorem }

On standard Borel spaces, the canonical conditional-law state
\[
Z_t
=
\mathbb P(Y_t^+\in\cdot\mid\mathbf H_t)
\]
is sufficient and minimal:
\[
Y_t^+\indep\mathbf H_t\mid Z_t,
\]
and every sufficient statistic \(S_t=s(\mathbf H_t)\) satisfies
\[
Z_t=\phi(S_t)
\qquad\text{a.s.}
\]
Moreover, a structural candidate \(W_t\) is exactly operationally aligned if
and only if
\[
Z_t=\phi_W(W_t)
\qquad\text{a.s.}
\]
The countable intervention-indexed version holds for
\(Z_t^{\mathcal I}\).

\subtopic{\(T_{13}\): Operational relationality theorem }

For a finite target, exact alignment, a positive node-level conditional
information gap, a positive robust local-predictive-matched null gap, a feasible
workspace intervention effect, and positive matched specificity jointly imply
that:
\[
\begin{aligned}
&\text{no single-node statistic is sufficient},\\
&k_0\geq2,\\
&\text{no single node reconstructs the aligned workspace},\\
&\text{matched local nulls preserve all node risks but lose joint capacity},\\
&\text{and the workspace coordinate is causally efficacious at the tested scale.}
\end{aligned}
\]
If \(\eta_{\mathrm{loop}}(\tau)>0\), the conclusion additionally includes
loop-mediated reentry.

\subtopic{\(T_{14}\): Replacement invariance theorem }

For a unitary realization correspondence with discrepancy defect
\[
\delta_D=\|D'J_V-J_ED\|,
\]
set
\[
\eta_L=2\|D\|\delta_D+\delta_D^2.
\]
If \(\eta_L<g/2\), then the workspace projectors obey
\[
\boxed{
q_{\mathrm{str}}
\leq
\frac{\eta_L}{g-\eta_L}
\leq
\frac{2\eta_L}{g}.
}
\]
A canonical unitary correspondence between the workspace subspaces exists,
and its operation defect is bounded by the projector drift, the local
intertwining defect, and prior operation leakage. Coalition frame bounds also
persist under sufficiently small reconstructive transport error. With declared
causal-kernel and reconstructive-handoff defects,
\[
\boxed{
d_{\mathrm{id}}
\leq
 a\,q_{\mathrm{str}}
+b\,\delta_{\mathrm{op}}
+c\,\delta_{\mathrm{causal}}
+d\,\delta_{\mathrm{rec}}.
}
\]
Under bounded norms and a uniform spectral gap, this gives the schematic
estimate
\[
d_{\mathrm{id}}
\lesssim
C\frac{\delta_{\mathrm{replacement}}}{g}.
\]
The causal and reconstructive terms are necessary assumptions rather than
consequences of structural similarity.


\subtopic{\(T_{15}\): Transfer identification theorem }

Under the linear transfer model
\[
Y=H_*X+\xi,
\qquad
\mathbb E[\xi X^{\mathsf T}]=0,
\qquad
\Sigma_{XX}\succeq\kappa I,
\]
the unique population least-squares transfer is
\[
\boxed{H_*=\Sigma_{YX}\Sigma_{XX}^{-1},}
\]
and
\[
\mathcal R(H)-\mathcal R(H_*)
\geq
\kappa\|H-H_*\|_F^2.
\]
For empirical ordinary least squares,
\[
\|\widehat H-H_*\|_F
\leq
\frac{\|\widehat M\|_F}{\widehat\kappa},
\]
with the stated ridge extension. The identified transfer is gauge covariant
under independent invertible changes of source and target coordinates.

\subtopic{\(T_{16}\): Uniform structure-selection theorem }

For a finite candidate family with unique population optimum \(\theta_*\),
margin \(\Delta_J>0\), and held-out uniform score error at most \(\varepsilon\),
\[
\boxed{
J(\widehat\theta)-J(\theta_*)\leq2\varepsilon.
}
\]
If \(2\varepsilon<\Delta_J\), then \(\widehat\theta=\theta_*\). The theorem
applies to node membership, edge selection, gauge-normalized transfer models,
and finite boundary families.

\subtopic{\(T_{17}\): Stable boundary recovery theorem }

For a finite admissible boundary family with unique maximizer \(S_t^*\), margin
\(\Delta_t^B\), and uniform held-out score error \(\varepsilon_t\),
\[
\mathcal B_t(S_t^*)-
\mathcal B_t(\widehat S_t)
\leq2\varepsilon_t,
\]
and exact recovery holds when \(2\varepsilon_t<\Delta_t^B\). If the population
score drifts uniformly by at most \(\delta_t\), then
\[
2\delta_t<\Delta_t^B
\Longrightarrow
S_{t+1}^*=S_t^*.
\]

\subtopic{\(T_{18}\): Model-specific formation transition theorem }

For
\[
m_{t+1}=\tanh(\lambda m_t),
\]
there is a unique zero fixed point for \(0\leq\lambda\leq1\). For
\(\lambda>1\), the zero state is unstable and exactly two stable nonzero fixed
points appear. Their amplitude satisfies
\[
\boxed{
m_\lambda
=
\sqrt{\frac{3(\lambda-1)}{\lambda^3}}
+
O\bigl((\lambda-1)^{3/2}\bigr).
}
\]
This is an exact structural benchmark, not a universal law for complete
relational workspaces.

\subtopic{\(T_{19}\): Cue-controlled goal-retention theorem }

For matched future command and disturbance trajectories under
\[
G_{s+1}=R_sG_s+B_sC_{s+1}+\zeta_{s+1},
\]
two goal interventions satisfy
\[
\boxed{
G_{t+\tau}-G_{t+\tau}'
=
R_{t+\tau:t}(g-g').
}
\]
If
\[
m_s\leq\sigma_{\min}(R_s),
\qquad
\|R_s\|\leq M_s,
\]
then
\[
\boxed{
\left(\prod_s m_s\right)\|g-g'\|
\leq
\|G_{t+\tau}-G_{t+\tau}'\|
\leq
\left(\prod_s M_s\right)\|g-g'\|.
}
\]
A linear action readout gives a corresponding lower bound on the induced
action-mean separation. The result isolates internal retention from repeated
external goal rewriting.

\subtopic{\(T_{20}\): Approximate counterfactual-policy theorem }

If the interventional action-value estimator satisfies
\[
\sup_a|\widehat Q_t(a)-Q_t(a)|\leq\varepsilon_Q
\]
and the implemented policy has optimization error at most
\(\delta_{\mathrm{opt}}\) relative to \(\widehat Q_t\), then
\[
\boxed{
\operatorname{Reg}_t(\widehat\pi_t)
\leq
2\varepsilon_Q+\delta_{\mathrm{opt}}.
}
\]
For deterministic greedy action, the bound is \(2\varepsilon_Q\). A unique
true optimum with value gap greater than \(2\varepsilon_Q\) is recovered
exactly. The guarantee is relative to the declared utility, feasible action
support, causal identification, and horizon.

\subtopic{\(T_{21}\): Repair-selection margin theorem }

For a finite repair set and a repair-value estimator with uniform error
\(\varepsilon_R\), the selected repair \(\widehat r\) satisfies
\[
\boxed{
V(r^*)-V(\widehat r)
\leq
2\varepsilon_R.
}
\]
If the optimal repair reaches reserve \(\rho+2\varepsilon_R\), then the selected
repair reaches reserve \(\rho\). A unique optimal repair whose value gap
exceeds \(2\varepsilon_R\) is selected exactly.

\subtopic{\(T_{22}\): Closed-loop self-maintenance theorem }

If matched no-repair deficit obeys
\[
e_{t+1}^{(0)}\leq\alpha e_t+d_t
\]
and the implemented repair has measured causal effect
\[
e_{t+1}^{(0)}-e_{t+1}\geq\kappa e_t-\epsilon_t,
\]
set \(\beta=\alpha-\kappa\). Under
\[
0\leq\beta<1,
\qquad
 d_t\leq\overline d,
\qquad
 \epsilon_t\leq\overline\epsilon,
\]
\[
\boxed{
e_t
\leq
\beta^te_0
+
\frac{(\overline d+\overline\epsilon)(1-\beta^t)}{1-\beta}.
}
\]
Hence
\[
\limsup_{t\to\infty}e_t
\leq
\frac{\overline d+\overline\epsilon}{1-\beta}.
\]
A corresponding margin condition makes the declared bounded-deficit viability
set forward invariant.

\subtopic{\(T_{23}\): Workspace-agency synthesis theorem }

Operational relationality, cue-controlled and behaviorally operative goals,
calibrated counterfactual values, feasible value-to-action and action-to-outcome
interventions, causal closed-loop maintenance, distributed control-state
sufficiency, and matched local policy nulls jointly imply:
\[
\begin{aligned}
&\operatorname{Reg}_t\leq2\varepsilon_Q+\delta_{\mathrm{opt}},\\
&\limsup e_t\leq
 \frac{\overline d+\overline\epsilon}{1-\beta},\\
&k_{\mathrm{ctl},\varepsilon}\geq2,
\end{aligned}
\]
together with causal mediation of goal-dependent action and repair through the
implemented internal channels. This is a sufficient, domain-relative agency
certificate. It does not imply consciousness, personhood, or unrestricted
autonomy.

\subtopic{\(T_{24}\): Gatewise certification-validity theorem }

Suppose a certificate requires every gate \(G_1,\ldots,G_m\) to pass, and the
predeclared test for gate \(j\) has false-pass probability at most \(\alpha_j\)
whenever that required property is absent. Under the global null that at least
one required property is absent,
\[
\boxed{
\mathbb P(\text{complete certificate passes})
\leq
\max_j\alpha_j.
}
\]
No independence assumption among gate tests is required. The guarantee fails
under adaptive gate deletion, selective reporting, or reuse of the
certification split.


\topic{70}{Decomposition, Specification, and Identification Results}

\subtopic{\(T_{25}\): Faithful coarse-graining theorem}
For an exact partition into internally coherent macro-components,
\[
\boxed{
\ker D_H=J_\Pi\ker D_\Pi.
}
\]
Thus faithful aggregation preserves the exact shared workspace and its
dimension while changing the displayed node resolution.

\subtopic{\(T_{26}\): Bayes-risk transport under aggregation}
For an information-faithful coarsening map \(q\),
\[
\boxed{
\Risk_{\mathcal D_c}(S)
=
\Risk_{\mathcal D_f}\bigl(q^{-1}(S)\bigr).
}
\]
Predictive content is preserved, although the visible coalition count may
contract. Primitive-carrier support cost is therefore preferred when claims
must be compared across resolutions.

\subtopic{\(T_{27}\): Decomposition-transport robustness theorem}
If every required gate has reference margin larger than its worst admissible
decomposition-transport defect, then the complete certificate passes
throughout \(\mathfrak D\). When the mismatch term vanishes, the certificate
radius is bounded below by the smallest gate margin divided by its transport
sensitivity.

\subtopic{\(T_{28}\): Specification-transport robustness theorem}
For a reference specification \(\varsigma_0\), if
\[
m_q(W;\varsigma_0)
>
\sup_{\varsigma\in\mathfrak S}
\left[
L_qd_{\mathfrak S}(\varsigma,\varsigma_0)+b_q(\varsigma)
\right]
\]
for every gate, then the transported workspace passes every admissible
specification. If \(b_q=0\),
\[
r_{\mathrm{spec}}(W;\varsigma_0)
\geq
\min_{q:L_q>0}\frac{m_q(W;\varsigma_0)}{L_q}.
\]

\subtopic{\(T_{29}\): Experimental non-identifiability theorem}
If two realizations induce the same law under every declared observation and
intervention but have different certificate values, no decision rule based on
that experiment can be uniformly correct. For every such rule, the two error
probabilities sum to one, so the minimax error is at least one half.

\subtopic{\(T_{30}\): Joint specification-identification theorem}
Let \(r_{\mathfrak S}(\mathfrak R)\) be the worst gate margin across the
specification family. If its infimum over the identified realization set is
positive, the robust certificate is identified as passing. If its supremum is
nonpositive, the certificate is identified as failing. Compatible realizations
on opposite sides of zero establish non-identifiability.

\subtopic{\(T_{31}\): Uniform finite-sample specification theorem}
If all gate estimators are uniformly within their error radii across
\(\mathfrak S\), then requiring every empirical lower envelope to be positive
controls the probability of a false robust certificate by the failure
probability of the uniform event.

\subtopic{\(T_{32}\): Independent selection and certification theorem}
A specification chosen on data independent of the certification split may be
evaluated by its fixed-specification test without inflating the false-pass
bound. Selection and certification on the same outcomes require simultaneous
or post-selection correction.
